\documentclass[english, parskip=full]{scrartcl}
\usepackage[utf8]{inputenc}  % Kodierung der Datei
\usepackage[english]{babel}  % Sprache auf Deutsch 
\usepackage{color}
\usepackage{graphicx, gensymb}
\usepackage{amsfonts, cancel, amssymb}
\usepackage{amsmath, amsthm, array, esvect, esint}
\usepackage{booktabs, multirow}  % schönere Tabellen
\usepackage{caption}  % Anpassung der Captions
\usepackage{scrlayer-scrpage}    % Manipulation des Seitenstils
\pagestyle{scrheadings}  % Stil für die Seite setzen
\automark{section}  % Abschnittsnamen als Seitenbeschriftung 
\setheadsepline{.5pt}    % Kopzeile durch Linie abtrennen
\usepackage[hidelinks]{hyperref}  % Links und weitere PDF-Features
\usepackage{typearea, tikz, pgffor, verbatim}
\usetikzlibrary{calc, quotes}
\usepackage[cache=false, outputdir=build]{minted}

\definecolor{bg}{rgb}{0.9,0.9,0.9}
\newcommand\numberthis{\addtocounter{equation}{1}\tag{\theequation}}
\newcommand{\bra}{}
\newcommand{\kom}{}
\newcommand{\brak}{}
\newcommand{\braket}{}
\newcommand{\intx}{}
\newcommand{\intr}{}
\newcommand{\kla}{}
\newcommand{\klam}{}
\newcommand{\klammer}{}
\newcommand{\oper}{}
\newcommand{\tecxt}{}
\newcommand{\Bra}{}
\newcommand{\Ket}{}
\def\I{\text{i}}

\newcommand*{\titel}{02 Lorentz group}
\newcommand*{\autor}{Joachim Schwardt}

\hypersetup{pdfauthor={\autor}, pdftitle={\titel}}

\subject{Quantum theory 2}
\title{\titel}
\author{\autor}
\date{\today}

\begin{document}

%\begin{titlepage}
\maketitle  % Titel setzen
%\tableofcontents  % Inhaltsverzeichnis setzen
%\end{titlepage}

\section{Questions}
Three rotation generators and three boost generators.\\
Let $x'^\mu=x^\mu +\omega^\mu_{\ \nu}x^\nu$ with $\omega^{\mu\nu}=-\omega^{\nu\mu}$. Then the scalar product is invariant:
\begin{align*}
x'^\mu x'_\mu&=x'^\mu g_{\mu\nu}x'^\nu=x^\mu g_{\mu\nu}x^\nu +\omega^\mu_{\ \rho}x^\rho g_{\mu\nu} x^\nu + x^\mu g_{\mu\nu}\omega^\nu_{\ \lambda}x^\lambda +\mathcal{O}(\omega^2)=\\
&=x^\mu x_\mu+x^\mu \kla\left( {\omega_{\nu\mu}+\omega_{\mu\nu}} \right)x^\nu +\mathcal{O}(\omega^2)\simeq x^\mu x_\mu 
\end{align*}
Left handed and right handed 2-spinors have different generators and lead to problems with defining the parity operator. Need Dirac 4-spinor to define this.\\
Let $\psi,\chi$ be Dirac spinors with $\overline{\psi}=\psi^\dagger\gamma^0$. Then their product is invariant.
\section{Generators of the Lorentz group}
We want to show that $\omega^\mu_{\ \nu}=-\frac{\I}{2}\omega^{\rho\sigma}(L_{\rho\sigma})^\mu_{\ \nu}$ for two specific cases.\\
(a): Rotations around the $z$-axis with $\omega^{12}=-\omega^{21}=\epsilon$:
\begin{align*}
-\frac{\I}{2}\omega^{12}(L_{12})^\mu_{\ \nu}-\frac{\I}{2}\omega^{21}(L_{21})^\mu_{\ \nu}=-\I\epsilon (l_z)^\mu_{\ \nu}=\begin{pmatrix}
0&0&0&0\\ 0&0&-\epsilon&0\\ 0&\epsilon&0&0\\ 0&0&0&0
\end{pmatrix}=g_{\lambda\nu}\omega^{\mu\lambda}=\omega^\mu_{\ \nu}
\end{align*}
(b): Boosts along the $x$-axis with $\omega^{10}=-\omega^{01}=\beta$:
\begin{align*}
-\frac{\I}{2}\omega^{10}(L_{10})^\mu_{\ \nu}-\frac{\I}{2}\omega^{01}(L_{01})^\mu_{\ \nu}=-\I\beta (k_x)^\mu_{\ \nu}=\begin{pmatrix}
0&\beta&0&0\\ \beta&0&0&0\\ 0&0&0&0\\ 0&0&0&0
\end{pmatrix}=g_{\lambda\nu}\omega^{\mu\lambda}=\omega^\mu_{\ \nu}
\end{align*}
\section{Lorentz transformation of spinors}
Let $\psi'=S(\Lambda)\psi$ with $S(\Lambda)=1-\frac{\I}{2}\omega^{\rho\sigma}S_{\rho\sigma}$. Then the following product is Lorentz invariant:
\begin{align*}
\overline{\psi}'\psi'&=(S\psi)^\dagger\gamma^0S\psi=\psi^\dagger S^\dagger\gamma^0S\psi=\overline{\psi}\psi
\end{align*}
For this, we need to show that $S^{-1}=\gamma^0S^\dagger\gamma^0$, or equivalently:
\begin{align*}
\gamma^0S^\dagger\gamma^0S&=\gamma^0\kla\left( {1+\frac{\I}{2}(\omega^{\rho\sigma})^\dagger S_{\rho\sigma}^\dagger} \right)\gamma^0S=\gamma^0\kla\left( {1+\frac{\I}{2}\omega^{\rho\sigma}\gamma^0S_{\rho\sigma}\gamma^0} \right)\gamma^0S=\\
&=\kla\left( {1+\frac{\I}{2}\omega^{\rho\sigma}S_{\rho\sigma}} \right)\kla\left( {1-\frac{\I}{2}\omega^{\rho\sigma}S_{\rho\sigma}} \right)=1+\mathcal{O}(\omega^2)
\end{align*}
We can also proof the following equality for infinitesimal transformations using $\klam\left[ {A,BC} \right]=B\klammer\left\{{A,C} \right\}-\klammer\left\{{A,B} \right\}C$:
\begin{align*}
S^{-1}\gamma^\mu S&=\kla\left( {1+\frac{\I}{2}\omega^{\rho\sigma}S_{\rho\sigma}} \right)\gamma^\mu\kla\left( {1-\frac{\I}{2}\omega^{\rho\sigma}S_{\rho\sigma}} \right)=\gamma^\mu +\frac{\I}{2}\omega_{\rho\sigma}S^{\rho\sigma}\gamma^\mu - \gamma^\mu \frac{\I}{2}\omega_{\rho\sigma}S^{\rho\sigma}=\\
&=\gamma^\mu -\frac{\omega_{\rho\sigma}}{8}\klam\left[ {\gamma^\mu, \klam\left[ {\gamma^\rho,\gamma^\sigma} \right]} \right]=\\
&=\gamma^\mu -\frac{\omega_{\rho\sigma}}{8}\kla\left( {\gamma^\rho\klammer\left\{{\gamma^\mu,\gamma^\sigma} \right\}-\klammer\left\{{\gamma^\mu,\gamma^\rho} \right\}\gamma^\sigma - \gamma^\sigma\klammer\left\{{\gamma^\mu,\gamma^\rho} \right\}+\klammer\left\{{\gamma^\mu,\gamma^\sigma} \right\}\gamma^\rho} \right)=\\
&=\gamma^\mu + \frac{\omega_{\rho\sigma}}{8}\cdot 4\gamma^\sigma\klammer\left\{{\gamma^\mu,\gamma^\rho} \right\}=\gamma^\mu +\omega_{\rho\sigma}\gamma^\sigma g^{\mu\rho}=\Lambda^\mu_{\ \nu}\gamma^\nu
\end{align*}
The quantity $\overline{\psi}\gamma^\mu\psi$ transforms as a 4-vector:
\begin{align*}
\overline{\psi}'\gamma^\mu\psi'&=\overline{\psi}S^{-1}\gamma^\mu S\psi=\overline{\psi}\Lambda^\mu_{\ \nu}\gamma^\nu\psi=\Lambda^\mu_{\ \nu}\overline{\psi}\gamma^\nu\psi
\end{align*}
\section{Commutator of Lorentz group operators}
Let $S_z=S_{12}$, $K_x=S_{10}$ and $K_y=S_{20}$. With the relation
\begin{align*}
S_{\mu\nu}=\frac{\I}{4}\klam\left[ {\gamma_\mu,\gamma_\nu} \right]=\frac{\I}{4}\klammer\left\{{\gamma_\mu,\gamma_\nu} \right\}-\frac{\I}{2}\gamma_\mu\gamma_\nu=\frac{\I}{2}\kla\left( {g_{\mu\nu}-\gamma_\nu\gamma_\mu} \right)\kom\overset{\substack{{\mu\neq\nu} \\ |}}{=}\frac{\I}{2}\gamma_\mu\gamma_\nu
\end{align*}
we can then compute the commutator 
\begin{align*}
\klam\left[ {S_z,K_x} \right]&=-\frac{1}{4}\klam\left[ {\gamma_2\gamma_1,\gamma_0\gamma_1} \right]=-\frac{1}{4}\klam\left[ {\gamma_2,\gamma_0\gamma_1} \right]\gamma_1-\frac{1}{4}\gamma_2\klam\left[ {\gamma_1,\gamma_0\gamma_1} \right]=\\
&=\I\gamma_0S_{12}\gamma_1+\I K_y+\I\gamma_2S_{10}\gamma_1=\I K_y+\frac{1}{2}\gamma_0\gamma_1\gamma_2\gamma_1+\frac{1}{2}\gamma_2\gamma_1\gamma_0\gamma_1=\\
&=\I K_y +\frac{1}{2}\begin{pmatrix}
-\sigma^1&0 \\ 0&\sigma^1
\end{pmatrix}\begin{pmatrix}
\I\sigma^3&0 \\ 0&\I\sigma^3
\end{pmatrix}+\frac{1}{2}\begin{pmatrix}
\I\sigma^3&0 \\ 0&\I\sigma^3
\end{pmatrix}\begin{pmatrix}
-\sigma^1&0 \\ 0&\sigma^1
\end{pmatrix}=\I K_y
\end{align*}
or, do the same in explicit matrix formulation. The latter  turns out to be more elegant.
\section{Commutators of the Dirac Hamiltonian with angular momenta}
The Dirac equation is $(\I\gamma^\mu\partial_\mu-m)\psi=0$. Multiplying by $\gamma^0$ yields
\begin{align*}
0&=(\I\partial_0+\I\gamma^0\gamma^j\partial_j-\gamma^0m)\psi\\
\Rightarrow\ \ \ \I\partial_t\psi&=\begin{pmatrix}
I_2m&-\I\sigma^j\partial_j \\ -\I\sigma^j\partial_j & -I_2m
\end{pmatrix}\psi=\begin{pmatrix}
I_2m&\hat{\vec{\sigma}}\cdot\hat{\vec{p}} \\ \hat{\vec{\sigma}}\cdot\hat{\vec{p}} & -I_2m
\end{pmatrix} =:H_\tecxt\text{D}\psi,
\end{align*}
where
$$H_\tecxt\text{D}=\begin{pmatrix}
m&0&-\I\partial_z&-\I\partial_x-\partial_y\\
0&m&-\I\partial_x+\partial_y&\I\partial_z\\
-\I\partial_z&-\I\partial_x-\partial_y&-m&0\\
-\I\partial_x+\partial_y&\I\partial_z&0&-m
\end{pmatrix}$$
Let $S_{12}=\begin{pmatrix}
\sigma^3&0\\ 0&\sigma^3
\end{pmatrix} $ and $\hat{\vec{S}}=\begin{pmatrix}
0&\hat{\vec{\sigma}} \\ \hat{\vec{\sigma}} & 0
\end{pmatrix} $:
\begin{align*}
\klam\left[ {S_{12},H_\tecxt\text{D}} \right]&=\begin{pmatrix}
\sigma^3&0\\ 0&\sigma^3
\end{pmatrix}\begin{pmatrix}
I_2m&\hat{\vec{\sigma}}\cdot\hat{\vec{p}} \\ \hat{\vec{\sigma}}\cdot\hat{\vec{p}} & -I_2m
\end{pmatrix}-\begin{pmatrix}
I_2m&\hat{\vec{\sigma}}\cdot\hat{\vec{p}} \\ \hat{\vec{\sigma}}\cdot\hat{\vec{p}} & -I_2m
\end{pmatrix}\begin{pmatrix}
\sigma^3&0\\ 0&\sigma^3
\end{pmatrix}=\\
&=\begin{pmatrix}
0&\klam\left[ {\sigma^3,\sigma^j} \right]p^j \\ \klam\left[ {\sigma^3,\sigma^j} \right]p^j & 0
\end{pmatrix}=\begin{pmatrix}
0&\I\sigma^2p^1-\I\sigma^1p^2\\
\I\sigma^2p^1-\I\sigma^1p^2&0
\end{pmatrix}=\I (\hat{\vec{p}}\times\hat{\vec{S}})_3
\end{align*}
Let $L_{12}=\I(x_1\partial_2-x_2\partial_1)$. The only non-trivial entries of $\klam\left[ {L_{12},H_\tecxt\text{D}} \right]$ are FIXME:
\begin{align*}
\klam\left[ {L_{12},p_1\pm\I p_2} \right]&=\klam\left[ {\I x_1\partial_2,\I\partial_1} \right]\pm\klam\left[ {\I x_1\partial_2,\partial_2} \right]-\klam\left[ {\I x_2\partial_1,\I\partial_1} \right]\mp\klam\left[ {\I x_2\partial_1,\partial_2} \right]=\\
&=-\partial_2\mp\I\partial_1\\
\Rightarrow\ \ \ \klam\left[ {L_{12},H_\tecxt\text{D}} \right]&=\begin{pmatrix}
0&0&0&-\partial_2-\I\partial_1\\
0&0&-\partial_2+\I\partial_1&0\\
0&-\partial_2-\I\partial_1&0&0\\
-\partial_2+\I\partial_1&0&0&0
\end{pmatrix}
\end{align*}
FIXME: I must have lost an i somewhere...\\
As expected $\klam\left[ {J_{12},H_\tecxt\text{D}} \right]=0$ for $J_{12}=L_{12}+S_{12}$. Thus, the total angular momentum is a conserved quantity, whereas spin and orbital angular momentum are not necessarily time independent individually.
\end{document}